\section{Pastry}

Pastry est une infrastructure distribuée de localisation d’objets et de routage destinée aux systèmes pair-à-pair (P2P) à grande échelle. Elle repose sur une table de hachage distribuée (DHT), comparable à Chord, et fonctionne au-dessus d’Internet via un réseau overlay. Les nœuds logiques de Pastry ne correspondent pas nécessairement à la topologie physique du réseau IP, mais forment une structure virtuelle organisée selon l’espace des identifiants.

L’objectif fondamental de Pastry est d’acheminer efficacement un message associé à une clé donnée vers le nœud dont l’identifiant est numériquement le plus proche de cette clé. À l’instar de Chord, Pastry résout le problème du \textit{key-based routing}, mais adopte une approche différente fondée sur un routage par préfixes plutôt que sur une structure strictement circulaire.

\subsection{Architecture générale}

La conception de Pastry repose sur plusieurs principes structurants :

\subsubsection*{1. Décentralisation}
Aucun nœud ne possède de rôle central. Chaque participant maintient localement ses structures de routage, ce qui supprime tout point unique de défaillance.

\subsubsection*{2. Scalabilité logarithmique}
Le routage s’effectue en $O(\log N)$ sauts en moyenne, où $N$ est le nombre de nœuds du système, ce qui garantit le passage à grande échelle.

\subsubsection*{3. Auto-organisation}
Le protocole s’adapte dynamiquement aux arrivées et aux départs de nœuds (\textit{churn}) par la mise à jour progressive des structures de routage.

\subsubsection*{4. Localité réseau}
Lorsque plusieurs nœuds satisfont une contrainte de routage, Pastry privilégie celui qui est le plus proche selon une métrique réseau, réduisant ainsi la latence réelle des communications.

\subsection{Identifiants et espace de clés}

Chaque nœud Pastry possède un identifiant unique, appelé \textit{nodeId}, généralement codé sur 128 bits et dérivé d’une fonction de hachage cryptographique. De même, chaque message ou clé à router reçoit un identifiant correspondant dans le même espace.

L’espace d’identifiants est organisé en base $2^{b}$, où $b$ est un paramètre de configuration influençant la taille des tables de routage. Le routage s’effectue en recherchant le nœud actif dont le \textit{nodeId} est numériquement le plus proche de l’identifiant cible.

Contrairement à Chord, qui dispose d’un anneau strictement circulaire, Pastry utilise un routage par préfixes, où la distance est mesurée par la longueur du préfixe commun entre identifiants.

\subsection{Structures de routage}

Chaque nœud Pastry maintient trois structures principales : la \textit{Routing Table}, le \textit{Leaf Set} et le \textit{Neighbor Set}.

\subsubsection{1. Routing Table}

Organisée par préfixes d’identifiants, elle contient des pointeurs vers des nœuds partageant des préfixes communs avec le \textit{nodeId} courant. Cette structure permet de faire des sauts logarithmiques vers la destination, rapprochant le message à chaque étape.

Cette structure est analogue aux \textit{Finger Tables} de Chord, mais basée sur des préfixes plutôt que sur des puissances de 2.

\subsubsection{2. Leaf Set (ensemble des feuilles)}

Il regroupe les nœuds immédiatement supérieurs et inférieurs au \textit{nodeId}. Il assure le routage correct pour les clés proches et fournit des alternatives en cas de défaillance de certains nœuds.

\subsubsection{3. Neighbor Set (ensemble des voisins)}

Il contient des nœuds physiquement proches selon la topologie réseau, indépendamment des identifiants. Il est utilisé pour optimiser la latence des communications et améliorer la localité réelle du routage.

\subsection{Processus de routage}

Le routage dans Pastry s’effectue en rapprochant progressivement le message de son \textit{nodeId} cible grâce aux structures précédentes :

\begin{enumerate}
    \item \textbf{Vérification du Leaf Set} : si la clé cible se trouve dans l’intervalle couvert, le message est envoyé au nœud le plus proche numériquement.
    \item \textbf{Utilisation de la Routing Table} : si la clé n’est pas dans le Leaf Set, le message est transmis à un nœud partageant un préfixe plus long avec la clé.
    \item \textbf{Sélection par proximité physique} : si aucune entrée n’est disponible ou valide, un nœud du Leaf Set numériquement plus proche est choisi, en prenant en compte la distance réseau.
\end{enumerate}

À chaque étape, le message progresse vers un nœud plus proche de la destination, ce qui garantit la convergence du routage en $O(\log N)$ sauts en moyenne pour $N$ nœuds actifs.

\subsection{Algorithmes des nœuds et maintenance}

Chaque nœud Pastry exécute des algorithmes pour intégrer de nouveaux nœuds, maintenir les structures de routage et gérer les défaillances.

\begin{enumerate}
    \item \textbf{Intégration d’un nouveau nœud (\textit{join})}
    \begin{itemize}
        \item Un nœud entrant contacte un nœud existant pour s’intégrer au réseau.
        \item Les structures de routage (\textit{Routing Table}, \textit{Leaf Set}, \textit{Neighbor Set}) sont initialisées progressivement grâce aux informations des nœuds voisins.
    \end{itemize}

    \item \textbf{Maintenance dynamique (\textit{stabilize} et mise à jour du routage)}
    \begin{itemize}
        \item Les entrées de la \textit{Routing Table} sont périodiquement mises à jour pour refléter les départs ou arrivées de nœuds.
        \item Le \textit{Leaf Set} et le \textit{Neighbor Set} sont également ajustés pour maintenir la cohérence et la résilience locale.
    \end{itemize}

    \item \textbf{Tolérance aux pannes}
    \begin{itemize}
        \item En cas de défaillance d’un nœud, les messages sont redirigés vers le nœud le plus proche dans le \textit{Leaf Set} ou via la \textit{Routing Table}.
        \item Cette redondance garantit que le routage reste correct et efficace même avec un nombre important de nœuds défaillants.
    \end{itemize}
\end{enumerate}
